\documentclass[11pt]{article}
\usepackage[margin=1in]{geometry}
\usepackage{amsmath,amssymb,amsfonts,amsthm,mathtools,bm}
\usepackage[T1]{fontenc}
\usepackage{microtype}
\newtheorem{theorem}{Theorem}
\newtheorem{lemma}{Lemma}
\newtheorem{cor}{Corollary}
\newtheorem{defn}{Definition}
\newtheorem{principle}{Principle}
\newtheorem{statement}{Statement}
\theoremstyle{remark}
\newtheorem{statementr}{Statement}

\title{The Set-Indexed Tensor Product}
\author{James Taylor}
\date{Working draft derived from the 2003 thesis}
\begin{document}
\maketitle

\section{The set-indexed tensor product} \label{s:setten} \label{a:tensor}


Let $\ensuremath{T}$ be a finite set. For each element $\ensuremath{a} \in
\ensuremath{T}$, we associate 
a vector space $\ensuremath{W}_{\ensuremath{a}}$ to that element. We define the
\ensuremath{T}-indexed Cartesian product 
of these spaces to be 
\begin{equation*}
\ensuremath{\mathop{\times}_{\ensuremath{a} \in \ensuremath{T}}
\ensuremath{W}_{\ensuremath{a}} } \ensuremath{:=} \ensuremath{\{f | f: \ensuremath{T} \to \bigcup_{\ensuremath{a}
\in \ensuremath{W}}{\ensuremath{W}_{\ensuremath{a}} \text{ such that } f(\ensuremath{a}) \in
\ensuremath{W}_\ensuremath{a} \} }}.\footnote{We could also say that the Cartesian
product is the set of sections of 
the bundle over the discrete space $T$ whose fiber at $\ensuremath{a}$ is
$\ensuremath{W}_{\ensuremath{a}}$.  }
\end{equation*} 
Endowed with the usual vector space structure, this is the 
\ensuremath{T}-indexed direct
sum of the spaces.\footnote{
As an example, the tangent space of \ensuremath{{}^N\ensuremath{\mathbb{R}^{\ensuremath{3}}}}\ is the \ensuremath{q}-indexed direct sum of
the tangent 
spaces in physical space.
} 
We define the \ensuremath{T}-indexed tensor product \ensuremath{\bigotimes_{\ensuremath{a} \in \ensuremath{T}} \ensuremath{W}_{\ensuremath{a}} }\ to be the set of
multilinear
functionals defined on \ensuremath{\mathop{\times}_{\ensuremath{a} \in \ensuremath{T}}
\ensuremath{\ensuremath{W}^*}_{\ensuremath{a}} }, the \ensuremath{T}-indexed Cartesian product of the
dual spaces.    

If, as will be our case, the vector spaces $\ensuremath{W}_{\ensuremath{a}}$ are the
same vector space, say \ensuremath{W}, then simplifications
occur. Notationally, \ensuremath{{\ensuremath{W}}^{\ensuremath{T}}}\ denotes the Cartesian
product; \ensuremath{{\ensuremath{W}}^{{\otimes}{\ensuremath{T}}}}\  denotes the tensor product. As for the definition of
the Cartesian product, the
union can be replaced with the vector space itself. The
Cartesian product becomes
\begin{equation*}
\ensuremath{{\ensuremath{W}}^{\ensuremath{T}}} = \{f| f: \ensuremath{T} \to \ensuremath{W}\}.
\end{equation*}
From this point on, we will assume that we have only one vector
space. The more general case does not substantially change any
statements below, but it is irrelevant to the purpose of this paper. 

Let us consider a different approach to \ensuremath{{\ensuremath{W}}^{{\otimes}{\ensuremath{T}}}}. We will use the
standard equivalence class approach to tensors. We start by
introducing some notation. For $f \in \ensuremath{{\ensuremath{W}}^{\ensuremath{T}}}$, we write
\ensuremath{\mathop{\otimes}_{\ensuremath{a} \in \ensuremath{T}} \ensuremath{w}_{\ensuremath{a}} }, where $\ensuremath{w}_{\ensuremath{a}} \ensuremath{:=} f(\ensuremath{a})$. 
Given a decomposition of \ensuremath{T}\ into the subsets $\ensuremath{T}_{\alpha}$, $\alpha
\in A$, for some $A$, 
then we
certainly take as true 
 \[\mathop{\otimes}_{\alpha \in A} (\mathop{\otimes}_{\ensuremath{a} \in
\ensuremath{T}_{\alpha}} \ensuremath{w}_{\ensuremath{a}})
= \ensuremath{\mathop{\otimes}_{\ensuremath{a} \in \ensuremath{T}} \ensuremath{w}_{\ensuremath{a}} } = f.
\]
A special case of the above, and one with a special abuse of notation,
is when $A$ is a two element subset and one of the sets in the
decomposition has only one element, say $\ensuremath{a}'$. Then we would write 
\ensuremath{(\mathop{\otimes}_{ \substack{\ensuremath{a}
 \in \ensuremath{T} \\
 \ensuremath{a} \neq \ensuremath{a}'}} \ensuremath{w}_\ensuremath{a}) \otimes {{w_{\ensuremath{a}'}} }} which is equal to \ensuremath{\mathop{\otimes}_{\ensuremath{a} \in \ensuremath{T}} \ensuremath{w}_{\ensuremath{a}} }\ as well as
$w_{\ensuremath{a}'} \otimes (\mathop{\otimes}_{\substack{\ensuremath{a} \in
\ensuremath{T} \\ \ensuremath{a} \neq \ensuremath{a}'}} \ensuremath{w}_{\ensuremath{a}}) $.
We can then view the 
\ensuremath{T}-indexed tensor product as the equivalence classes of all formal
linear combinations of the elements in the set \ensuremath{{\ensuremath{W}}^{\ensuremath{T}}}\ under the
relations 
\begin{align*}
c(\ensuremath{\mathop{\otimes}_{\ensuremath{a} \in \ensuremath{T}} \ensuremath{w}_{\ensuremath{a}} }) &= \ensuremath{(\mathop{\otimes}_{ \substack{\ensuremath{a}
 \in \ensuremath{T} \\
 \ensuremath{a} \neq \ensuremath{a}'}} \ensuremath{w}_\ensuremath{a}) \otimes {{c\ensuremath{w}_{\ensuremath{a}'}} }} 
\\
\ensuremath{(\mathop{\otimes}_{ \substack{\ensuremath{a}
 \in \ensuremath{T} \\
 \ensuremath{a} \neq \ensuremath{a}'}} \ensuremath{w}_\ensuremath{a}) \otimes {{x_{\ensuremath{a}'}} }} + \ensuremath{(\mathop{\otimes}_{ \substack{\ensuremath{a}
 \in \ensuremath{T} \\
 \ensuremath{a} \neq \ensuremath{a}'}} \ensuremath{w}_\ensuremath{a}) \otimes {{y_{\ensuremath{a}'}} }} &= \ensuremath{(\mathop{\otimes}_{ \substack{\ensuremath{a}
 \in \ensuremath{T} \\
 \ensuremath{a} \neq \ensuremath{a}'}} \ensuremath{w}_\ensuremath{a}) \otimes {{(x
+y)_{\ensuremath{a}'}} }}.  
\end{align*}
These relations should hold for every $\ensuremath{a}' \in
\ensuremath{T}$, $c \in \ensuremath{\mathbb{C}}$ and
 $\ensuremath{w}_\ensuremath{a}$ $ \ensuremath{w}_{\ensuremath{a}'}$,  $x_{\ensuremath{a}'}$, $y_{\ensuremath{a}'} \in
\ensuremath{W}$.    
We should note the standard fact that
the multilinear functionals acting on \ensuremath{{\ensuremath{W}}^{\ensuremath{T}}}\ extend to
linear functionals acting on \ensuremath{{\ensuremath{W}}^{{\otimes}{\ensuremath{T}}}}; this is a third way of viewing \ensuremath{{\ensuremath{W}}^{{\otimes}{\ensuremath{T}}}}.
 
If \ensuremath{W}\ has an inner product, then there is a natural inner
product on \ensuremath{{\ensuremath{W}}^{{\otimes}{\ensuremath{T}}}}. Namely, for elements in \ensuremath{{\ensuremath{W}}^{\ensuremath{T}}}, we define the inner
product to be 
\begin{equation*}\ensuremath{({\ensuremath{\mathop{\otimes}_{\ensuremath{a} \in \ensuremath{T}} \ensuremath{w}_{\ensuremath{a}} }}, {\ensuremath{\mathop{\otimes}_{\ensuremath{a} \in
\ensuremath{T}} {{v}_{\ensuremath{a}} }}})} \ensuremath{:=}
\prod_{\ensuremath{a} 
\in \ensuremath{T}} \ensuremath{({\ensuremath{w}_\ensuremath{a}}, {v_{\ensuremath{a}}})}.
\end{equation*} 
We then extend this definition to \ensuremath{{\ensuremath{W}}^{{\otimes}{\ensuremath{T}}}}\ by using the appropriate
linearities.  

 If we have a linear operator $C$ acting on \ensuremath{W}, then
$C^{\ensuremath{a}}$ shall denote a linear operator acting on \ensuremath{{\ensuremath{W}}^{{\otimes}{\ensuremath{T}}}}\ such
that, when acting 
on product elements, it applies $C$ to the 
  $\ensuremath{a}^{th}$
factor while applying the identity to the other factors, i.e.\ 
\begin{equation*}
C^{\ensuremath{a}'} (\ensuremath{\mathop{\otimes}_{\ensuremath{a} \in \ensuremath{T}} \ensuremath{w}_{\ensuremath{a}} }) \ensuremath{:=} \ensuremath{(\mathop{\otimes}_{ \substack{\ensuremath{a}
 \in \ensuremath{T} \\
 \ensuremath{a} \neq \ensuremath{a}'}} \ensuremath{w}_\ensuremath{a}) \otimes {{C w_{\ensuremath{a}'}} }}
\end{equation*} 

We finish with a word about contractions. The set-indexed tensor product
is excellent for use with contractions. Indeed, given $w \in
\bigotimes_{a \in \alpha} E \otimes \bigotimes_{b \in \beta} E^{*}$,
then the contraction of $w$ over the index pair $(\tilde a, \tilde b)$
is denoted by $\ensuremath{\mathcal{C}}(\tilde a, \tilde b) w$ and is an element of
$\bigotimes_{a \in \{\alpha \setminus \{\tilde a\} \}} 
E \otimes \bigotimes_{b \in \{\beta \setminus \{\tilde b \}\} }
E^{*}$. This allows one to notationally keep track of which indices
were contracted. Using numbers as indices eliminates this
information. Furthermore, one can also see that a type of unit
analysis may be invoked. In other words, the expression $
\ensuremath{\mathcal{C}}(\tilde a, \tilde b) w + \ensuremath{\mathcal{C}}(a', {b'}) w$ would seem to make
no sense whereas 
$\ensuremath{\mathcal{C}}(\tilde a, \tilde b) \ensuremath{\mathcal{C}}(a',  b') w + 
\ensuremath{\mathcal{C}}(\tilde a,  b') \ensuremath{\mathcal{C}}( a', \tilde b) w$ would make sense. With
numbers, this would be hard to indicate. As such, the set-indexed tensor
product may be as
helpful in formulating meaningful expressions
 as dimensional analysis has historically been.

\clearpage
\bibliographystyle{abbrv}
\bibliography{references}
\end{document}
